\documentclass{article}
\usepackage[utf8]{inputenc}
\usepackage[T1]{fontenc}
\usepackage{geometry}
\geometry{
    a4paper,
    left=2.5cm,
    right=2.5cm,
    top=2.5cm,
    bottom=2.5cm
}
\usepackage{graphicx}
\usepackage{float}
\usepackage{hyperref}
\usepackage[english]{babel}
\usepackage{amsmath}
\usepackage{amssymb}
\usepackage{array}
\usepackage{booktabs}
\usepackage{caption}
\usepackage{subcaption}
\usepackage{enumitem}
\usepackage{fancyhdr}
\usepackage{titlesec}
\usepackage{url}
\usepackage{textcomp}
\usepackage{xcolor}
\usepackage{listings}
\usepackage{multirow}
\usepackage{longtable}
\usepackage{tabularx}
\usepackage{times}

% Set up page style
\pagestyle{fancy}
\fancyhf{}
\rfoot{\thepage}
\renewcommand{\headrulewidth}{0pt}

% Title formatting
\titleformat{\chapter}[display]{\normalfont\bfseries}{\chaptertitle\thechapter}{0pt}{\Huge}
\titlespacing*{\chapter}{0pt}{50pt}{40pt}

% Document info
\title{Template Document}
\author{Academic Support Agent}
\date{\today}

\begin{document}
\maketitle

%Her begynder indholdet af dokumentet.

%-------------------------------------------
%START HER!
%-------------------------------------------

%Forside
\begin{titlingpage}
\thispagestyle{empty}
\vspace*{3cm}
\centering
\Huge \bfseries \sffamily Skabelon til st rre projekter\\

\vspace*{3cm}
\includegraphics[width=7cm]{cirkler}

\vspace*{4cm}
\begin{flushright}
\large
Lavet af: Navn\\
Klasse: 2.x\\
Fag: Matematik\\
Dato: \today
\end{flushright}
\end{titlingpage}

%Indholdfortegnelse



\chapter{F rste kapitel}
Her er teksten til det f rste kapitel. Den kunne indeholde mange interessante ting.
\section{Afsnit}
Der er s rme ogs  et afsnit!

\chapter{Andet kapitel}\label{kapitel2}
Her er et andet kapitel, som altid starter p  en ny side. Her er det side \pageref{kapitel2}.

\chapter{Tabeller og figurer}

\section{Tabeller}
Hvis man vil lave en simpel tabel, kan man g re det p  f lgende m de
\begin{table}[htb]
\centering
\begin{tabular}{|l|r|}
\hline 
$n$ & Primtal mindre end eller lig $n$ \\ 
\hline 
1 & 0 \\ 
\hline 
2 & 1 \\ 
\hline 
3 & 2 \\ 
\hline 
4 & 2 \\ 
\hline 
\end{tabular}
\caption{Primtal mindre end eller lig et naturligt tal.}\label{tabel-primtal} 
\end{table}

Den tabel har nummeret \ref{tabel-primtal}. Tabellen er indlagt som en s kaldt \emph{float}, s  \LaTeX\ s tter den selv ind p  det sted i teksten, hvor den synes, der er bedst plads. S  beh ver man ikke bekymre sig om det.

\section{Figurer}
Man kan ogs  bruge floats til figurer. Det har vi gjort for figur \ref{fig-cirkler}.

\begin{figure}[htb]
\centering
\includegraphics[width=5cm]{cirkler}
\caption{To cirkler.}\label{fig-cirkler}
\end{figure}

Man kan l se mere i \cite[kapitel~6]{latex-bog}, hvor brugen af floats er beskrevet i langt st rre detalje.


%Litteraturliste
\begin{thebibliography}{1}
\bibitem{latex-bog}
Lars Madsen.
\textit{Introduktion til \LaTeX{}}.
2011.
\bibitem{anden-bog}
Jens Carstensen.
\textit{Komplekse tal}.
Systime,
3. udgave, 2. oplag, 1996.
\end{thebibliography}
%---------------------------------------------------------
%Dokumentet er slut.
%---------------------------------------------------------
\end{document}